\section{率失真：允许失真时的压缩极限}
\label{sec:07-Information-Theory/04-Source-Coding/05}

一张 \(1920\times1080\) 的 8 位 RGB 照片，原始像素约 6 MB。无损压缩能压到两三兆，JPEG 压到三百 KB 仍然看不出破绽。差出来的那一个数量级去哪了？

答案有点扎心：无损压缩里相当大一部分比特花在传感器噪声上。无损的定义要求每个像素一位不差地还原，而噪声恰恰是最不可压缩的东西——它没有结构，熵接近满值，谁也帮不了它。于是压缩器被迫为一堆没人要看的随机扰动认真买单。

一旦允许``差不多就行''，这笔钱就省下来了。问题是省到什么程度算过头。率失真理论要做的，就是把``允许多少误差''和``需要多少比特''之间的兑换率写成一个函数，让这件事不再靠试。

\subsection{把权衡写成一个下确界}

\begin{definition}
设信源 \(X\) 服从 \(p(x)\)，重建符号取值为 \(\hat x\)，失真度量 \(d(x,\hat x)\ge0\)。在平均失真约束 \(\E[d(X,\hat X)]\le D\) 之下，率失真函数定义为

\[
R(D)=\inf_{p(\hat x\given x):\ \E[d(X,\hat X)]\le D} I(X;\hat X).
\]
\end{definition}

要解释的是为什么被最小化的是互信息。\(I(X;\hat X)\) 度量重建保留了多少关于原始信号的信息，见\hyperref[sec:07-Information-Theory/02-Joint-Information/02]{``互信息是依赖，也是对不确定性的减少''}。而上一节已经把``信息''和``比特''钉在了一起：保留多少信息，就得传多少比特。所以在满足质量要求的所有编码方案里挑保留信息最少的那个，就是在挑码率最低的那个。约束管质量，目标管代价，两者被同一个量连起来。

这个定义的操作含义由 Shannon 率失真定理给出：对离散无记忆信源和满足条件的失真度量，长块编码可以在平均失真不超过 \(D\) 的前提下把码率压到任意接近 \(R(D)\)，而低于 \(R(D)\) 的码率无论怎么设计都达不到该失真水平。论证的骨架和上一节同源——先数一数在失真容限内``互相分不清''的重建点有多少，再把这个计数换成指数。更一般的信源需要相应的过程版本，单字母公式不能照搬。

\subsection{曲线的形状比闭式解重要}

\(R(D)\) 具体等于什么，大多数情况下算不出来。但它的形状是稳定的，而判断多半只需要形状。

\begin{center}
\begin{tikzpicture}[x=4.6cm,y=1.15cm,font=\small,>=stealth]
  \draw[->,black!55] (-0.03,0) -- (1.34,0) node[right,font=\footnotesize] {\(D/D_{\max}\)};
  \draw[->,black!55] (0,-0.06) -- (0,2.9) node[above,font=\footnotesize] {\(R(D)\)（比特）};
  \draw[black!20,dashed] (0,1) -- (1.25,1);
  \draw[black!20,dashed] (0,2) -- (1.25,2);
  \node[left,font=\footnotesize,black!60] at (-0.01,1) {\(1\)};
  \node[left,font=\footnotesize,black!60] at (-0.01,2) {\(2\)};
  \node[below,font=\footnotesize,black!60] at (0,-0.08) {\(0\)};
  \node[below,font=\footnotesize,black!60] at (1,-0.08) {\(1\)};
  \draw[blue!75,thick] plot[domain=0.032:1,samples=140] (\x,{-0.5*ln(\x)/ln(2)});
  \draw[blue!75,thick] (1,0) -- (1.25,0);
  \draw[red!70,thick,dashed] plot[domain=0.004:1,samples=140]
    (\x,{(-(0.25*ln(0.25)+0.75*ln(0.75))+(0.25*\x*ln(0.25*\x)+(1-0.25*\x)*ln(1-0.25*\x)))/ln(2)});
  \draw[red!70,thick,dashed] (1,0) -- (1.25,0);
  \fill[red!70] (0.004,0.811) circle (1.3pt);
  \draw[->,black!55] (1.12,0.58) -- (1.12,0.06);
  \node[above,font=\footnotesize,black!70,align=center] at (1.12,0.60) {\(D\ge D_{\max}\)：\\码率归零};
  \draw[blue!75,thick] (0.04,-0.66) -- (0.15,-0.66);
  \node[right,font=\footnotesize] at (0.16,-0.66) {Gaussian 信源、平方误差：左端发散};
  \draw[red!70,thick,dashed] (0.04,-1.04) -- (0.15,-1.04);
  \node[right,font=\footnotesize] at (0.16,-1.04) {二元信源、汉明失真：左端停在 \(H(X)\)};
\end{tikzpicture}

\emph{两条曲线都单调不增、都是凸的、都在 \(D=D_{\max}\) 处触零——那时最好的重建就是一个与数据无关的常数。差别全在左端：离散信源的无损极限是有限的 \(H(X)\)，连续信源没有有限的无损码率，因为精确写下一个实数要无穷多位。}
\end{center}

单调不增几乎是定义直接给的：放宽约束不会让下确界变大。凸性来自可以在两套方案之间做时间共享——一半时间用方案甲、一半用方案乙，失真和码率都取平均值，于是任何两点的连线都在曲线之上。凸性有个实用推论：在曲线上花码率的边际收益递减，前几个比特买到的质量提升最大，后面越来越贵。这就是为什么把一个视频从中等码率提到高码率，肉眼几乎看不出差别，而从极低码率提一点则效果显著。

右端点是全丢的位置。\(D_{\max}\) 定义为用一个最优常数去重建所能达到的失真，超过它再放宽约束也没有意义——码率已经是零。左端点则区分两类信源。离散信源在 \(D=0\) 处回到 \(H(X)\)，率失真理论此时退化成前几节的无损理论。连续信源的 \(R(0)\) 是无穷，这不是理论的缺陷，而是对``实数需要无限精度''这件事的正确记账。

Gaussian 信源在平方误差下有闭式：\(X\sim\mathcal N(0,\sigma^2)\) 时

\[
R(D)=\begin{cases}
\tfrac12\log_2\dfrac{\sigma^2}{D},&0<D<\sigma^2,\\[2pt]
0,&D\ge\sigma^2.
\end{cases}
\]

读法是把方差劈成两半：容许失真 \(D\) 相当于把方差为 \(D\) 的那一部分当噪声扔掉，只为方差 \(\sigma^2-D\) 的那部分付费。每砍一半失真，码率增加半个比特——这就是``\(6\) dB 换 \(1\) 比特''那条工程经验法则的来历。多维情形把这个动作对每个分量各做一次，方差小于门限的分量直接整个丢弃，方差大的按门限截断，形成反向注水的图景；主成分分析里``丢掉小特征值方向''的做法，在这里得到了码率意义上的解释。

\subsection{同一个权衡在表示学习里换了名字}

量化是率失真最直接的落地。把连续值映到有限个码点上，码点越多失真越小、索引越长，这正是曲线上的一次取点。神经网络的权重量化、向量量化的码本设计、扩散模型潜空间的压缩比选择，做的都是同一件事，区别只在失真度量换成了任务相关的东西。

再往前一步是信息瓶颈：把表示学习写成压缩 \(X\) 到 \(T\) 的同时保留关于 \(Y\) 的信息，目标形如 \(\min I(X;T)-\beta I(T;Y)\)，见\hyperref[sec:07-Information-Theory/06-Information-and-Learning/03]{``信息瓶颈与任务相关表示''}。外形和率失真完全一致——一个压缩代价，一个保真要求，一个换算系数——但两者约束的东西不同。率失真事先规定了逐样本的失真函数 \(d(x,\hat x)\) 并约束其均值；信息瓶颈约束的是表示保留了多少关于另一个变量的统计信息，没有逐样本的误差概念。某些设定下可以把后者改写成带特定失真的率失真问题，但那需要额外推导，不能仅凭形状相似就当成同一个定理。

即便如此，这个类比给出的判断是可靠的：``表示该保留多少信息''不是一个只能凭经验调的旋钮，它有一条可计算的权衡曲线，而曲线的凸性意味着极端两端都不划算。

\subsection{用之前要核对的几件事}

率失真定理和信源编码定理一样是渐近结论，块长趋于无穷才能贴到 \(R(D)\)。有限块长下实际码率更高，修正项与无损情形同族。

更容易出事的是失真度量的选择。\(R(D)\) 完全依赖 \(d\) 的定义：平方误差配 Gaussian 与部分听觉模型，汉明距离配二进制数据，图像的主观质量则和均方误差相去甚远——这就是 PSNR 高而观感差的情形为何常见，也是 SSIM、LPIPS 这类感知度量存在的理由。选错度量，理论极限依然成立，只是它保证的那个``质量''不是你要的质量。同样的道理在有损压缩的应用侧反复出现：说某种失真``可以接受''，必须连同任务一起说，医学影像和社交媒体缩略图不共享同一条曲线。

工程上没人直接求解率失真优化。JPEG 与 H.264 走的是变换加量化加熵编码的路线：先用 DCT 或小波把能量集中到少数系数上，再按视觉重要性分配量化步长，最后交给熵编码器。率失真给出理论极限，这套流程给出可实现的逼近，两者的差距就是编码标准迭代的空间。至于 \(R(D)\) 本身，除 Gaussian 等少数情形外没有显式解，需要用 Blahut--Arimoto 这类交替优化数值求解——它的结构和 EM 同源，交替更新条件分布与重建分布直到收敛。

允许失真解决了``必须精确还原''这个前提，但另一个前提一直没动过：从 Kraft 到算术编码到率失真，每一步都假设信源分布是已知的。真实的日志、源码和文本没有附带概率表，而估计一个分布本身就要花比特。下一节看看在完全不写下分布的情况下，压缩还能不能逼近熵率。
